% LTeX: language=de

\subsection{Übungsaufgabe}

\begin{Exercise}{Kurvendiskussion}{}

    Gegeben ist folgende gebrochen rationale Funktion: \qquad $f(x) = \dfrac{9-x^2}{1-x^2}$\\

    Bestimme folgende Eigenschaften der Funktion:
    \begin{itemize}
        \item Definitionsmenge $D_f$; Wertemenge $W_f$
        \item Nullstellen
        \item Verhalten am Rande der Definitionsmenge; Asymptoten
        \item Monotonie- und Krümmungsverhalten
    \end{itemize}

    Fertige zudem eine Zeichnung des Graphen $G_{f}$ an.

\end{Exercise}
\vspace{1ex}

\begin{solution}{}
    \begin{equation*}
        f(x) = \dfrac{-x^{\textbf{\textcolor{myr}{2}}} + 9}{-x^{\textbf{\textcolor{mydarkblue}{2}}} + 1} = \dfrac{z(x)}{n(x)}; \qquad \underbrace{z(x) = -x^2 + 9}_{\text{\textcolor{myg}{\text{\footnotesize{Zähler}}}}}; \qquad \underbrace{n(x) = -x^2 + 1}_{\text{\textcolor{myg}{\text{\footnotesize{Nenner}}}}}
    \end{equation*}
    Da $\text{\textbf{\textcolor{myr}{Zählergrad (ZG)}}} \geq \text{\textbf{\textcolor{mydarkblue}{Nennergrad (NG)}}}$ ist $f(x)$ eine \textbf{unecht gebrochen rationale Funktion}.
    \begin{itemize}[leftmargin=*]
        \item \begin{minipage}[t]{0.45\linewidth}
                Zählernullstellen:
                \begin{align*}
                    z(x) &= 0 \\
                    -x^2 + 9 &= 0 &&| + x^2\\
                    x^2 &= 9 &&| \sqrt[2]{\phantom{x}}\\
                    x_1 &= - 3 \text{\textit{(e)}} ; \qquad x_2 = 3 \text{\textit{(e)}} &&\\
                \end{align*}
                \vspace{1ex}
                \uuline{$N_1 \left(-3/0\right)$}; \quad \uuline{$N_2 \left(3/0\right)$}\\[2ex]
                $\Rightarrow$ \href{https://www.mathe-online.at/mathint/fun2/i_Nullstellen.html}{Nullstellen erster Ordnung} bei $x_1$ und $x_2$ jeweils mit Vorzeichenwechsel (VZW), da einfache Vielfachheit der Nst\\[1ex]
            \end{minipage}
            \hfill
            \begin{minipage}[t]{0.45\linewidth}
                Nennernullstellen:
                \begin{align*}
                    n(x) &= 0 \\
                    -x^2 + 1 &= 0 &&| + x^2\\
                    x^2 &= 1 &&| \sqrt[2]{\phantom{x}}\\
                    x_3 &= - 1 \text{\textit{(e)}} ; \qquad x_4 = 1 \text{\textit{(e)}} &&\\
                \end{align*}
                \vspace{1ex}
                \uuline{$\mathbb{D}_{f} = \mathbb{R} \backslash \{-1; 1 \}$}\\[2ex]
                $\Rightarrow$ \href{https://www.geogebra.org/m/Uz6x3YMq}{Polstellen erster Ordnung} bei $x_3$ und $x_4$ jeweils mit VZW, da ungerade Vielfachheit der Nst
            \end{minipage}
        \item Symmetrie:\qquad \begin{minipage}[t]{0.7\linewidth}
                \raisebox{-\height}{\begin{note}
                    Achsensymetrisch zur y-Achse: \qquad $f(-x) = f(x)$ \\
                    Punktsymetrisch zum Ursprung: \qquad $f(-x) = -f(x)$
                \end{note}}
                \vspace{1ex}
                \begin{flalign*}
                    \highlight{f(-x)} &= \dfrac{-(-x)^2 + 9}{-(-x)^2 + 1} = \dfrac{-x^2 + 9}{-x^2 + 1} \highlight{= f(x)}
                \end{flalign*}
                \vspace{1ex}
                $\Rightarrow$ $G_f$ ist achsensymmetrisch zur y-Achse
            \end{minipage}
        \item Asymptoten: \qquad \begin{minipage}[t]{0.45\linewidth}
            $\uuline{\textcolor{myorange}{v_1(x): \ x = -1}}$ \qquad $\uuline{\textcolor{myorange}{v_2(x): \ x = 1}}$\\[-3ex]
            \begin{flalign*}
            (&- x^2 &+ &\textcolor{mygr}{0 x} &+ &9) : (- x^2 + \textcolor{mygr}{0 x} + 1)
                \ = \ 1 + \smash[b]{\underbrace{\dfrac{8}{-x^2 + 1}}_{\text{\textcolor{myg}{\footnotesize Rest}}}} &\\[-0.5ex]
            \makebox[0pt][l]{\uline{\phantom{$xxxxxxxxxxxxxxx$}}}
            -(&- x^2 &+ &\textcolor{mygr}{0 x} &+ &1) &\\
            & \ \text{\textdiscount} & &\text{\textdiscount} & &8 &
            \end{flalign*}
            $\uuline{\textcolor{myorange}{h_1(x): \ y = 1}}$
        \end{minipage}
        \item Grenzverhalten: \hspace{-2cm}\begin{minipage}[t]{0.80\linewidth}
                \begin{equation*}
                    R(x) = \dfrac{8}{-x^{\textbf{\textcolor{mydarkblue}{2}}} + 1}
                \end{equation*}
                Da $\text{\textbf{\textcolor{myr}{ZG}}} < \text{\textbf{\textcolor{mydarkblue}{NG}}}$ ist $R(x)$ eine \textbf{echt gebrochen rationale Funktion}. Zur Bestimmung des Grenzverhaltens muss in der Regel zuerst die höchste Potenz ausgeklammert (\textbf{Normierung}) oder $z'(x)$ und $n'(x)$ gebildet (\textbf{Regel von L’Hospital}) werden, außer $\text{\textbf{\textcolor{mydarkblue}{NG}}} = 0 $. \\[-2.0ex]
                \begin{flalign*}
                    \lim\limits_{x \to -\infty} R(x) &= \dfrac{8}{-(\textcolor{mypurple}{-\infty})^2 + 1} = 0^- \quad &\rightarrow                     
                    \parbox[t]{5cm}{Annäherung von \textbf{unten} an \textcolor{myorange}{$h_1(x)$}}
                    &&\\[0.8ex]
                    \lim\limits_{x \to +\infty} R(x) &= \dfrac{8}{-(\textcolor{mypurple}{\infty})^2 + 1}= 0^- \quad &\rightarrow
                    \parbox[t]{5cm}{Annäherung von \textbf{unten} an \textcolor{myorange}{$h_1(x)$}}
                    &&\\[0.8ex]
                    \lim\limits_{x \to -1^-} R(x) &= \dfrac{8}{-(\textcolor{mypurple}{-1^-})^2 + 1} = -\infty \quad &\rightarrow 
                    \parbox[t]{5cm}{Linksseitige Annäherung an \textcolor{myorange}{$v_1(x)$} nach \textbf{unten}}
                    &&\\[0.8ex]
                    \lim\limits_{x \to -1^+} R(x) &= \dfrac{8}{-(\textcolor{mypurple}{-1^+})^2 + 1} = \infty \quad &\rightarrow 
                    \underbrace{\parbox[t]{5cm}{Rechtsseitige Annäherung an \textcolor{myorange}{$v_1(x)$} nach \textbf{oben}}}_{\text{\textcolor{myg}{\footnotesize \parbox[t]{5cm}{Bereits bekannt: Polstelle mit VZW}}}}
                    &&\\[0.8ex]
                    \lim\limits_{x \to 1^-} R(x) &= \dfrac{8}{-(\textcolor{mypurple}{1^-})^2 + 1} = \infty \quad &\rightarrow
                    \underbrace{\parbox[t]{5cm}{Linksseitige Annäherung an \textcolor{myorange}{$v_2(x)$} nach \textbf{oben}}}_{\text{\textcolor{myg}{\footnotesize \parbox[t]{5cm}{Bereits bekannt: \\ Achsensymmetrie zur y-Achse}}}}
                    &&\\[0.8ex]
                    \lim\limits_{x \to 1^+} R(x) &= \dfrac{8}{-(\textcolor{mypurple}{1^+})^2 + 1} = -\infty \quad &\rightarrow
                    \underbrace{\parbox[t]{5cm}{Rechtsseitige Annäherung an \textcolor{myorange}{$v_2(x)$} nach \textbf{unten}}}_{\text{\textcolor{myg}{\footnotesize \parbox[t]{5cm}{Achsensymmetrie zur y-Achse}}}}
                    &&
                \end{flalign*}
            \end{minipage}
        \item Monotonie- und Krümmungsverhalten: \\[1.5ex]
        \hspace{2cm} \begin{minipage}[t]{0.95\linewidth}
            \begin{note}
                Quotientenregel: \begin{minipage}[t]{0.8\linewidth}
                    \begin{align*}
                        f(x) &= \dfrac{\text{\textbf{\textcolor{myr}{Z}}}}{\text{\textbf{\textcolor{mydarkblue}{N}}}} \\[1.5ex]
                        f'(x) &= \dfrac{\text{\textbf{\textcolor{mydarkblue}{N}}} \cdot \text{\textbf{\textcolor{myr}{Z'}}} - \text{\textbf{\textcolor{myr}{Z}}} \cdot \text{\textbf{\textcolor{mydarkblue}{N'}}}}{(\text{\textbf{\textcolor{mydarkblue}{N}}})^2} = \dfrac{\text{\textbf{\textcolor{mydarkblue}{N}}}\text{\textbf{\textcolor{myr}{AZ}}} - \text{\textbf{\textcolor{myr}{Z}}}\text{\textbf{\textcolor{mydarkblue}{AN}}}}{(\text{\textbf{\textcolor{mydarkblue}{N}}})^2}
                    \end{align*}
                \end{minipage}
            \end{note}
            \begin{flalign*}
                f'(x) &= \dfrac{(-x^2 + 1) \cdot (-2x) - (-x^2 + 9)\cdot (-2x)}{(-x^2 + 1)^2} &&\\
                &= \dfrac{2x^3 - 2x - 2x^3 + 18x}{(-x^2 + 1)^2} &&\\
                &= \dfrac{16x}{(-x^2 + 1)^2}
            \end{flalign*}
            \begin{flalign*}
                f''(x) &= \dfrac{(-x^2 + 1)^2 \cdot 16 - 16x \cdot 2(-x^2 + 1)(-2x)}{(-x^2 + 1)^4} &&\\
                &= \dfrac{16(-x^2 + 1)^2 - 64x^2(-x^2 + 1)}{(-x^2 + 1)^4} &&\\
                &= \dfrac{16(-x^4 + 2x^2 - 1) + 64x^4 - 64x^2}{(-x^2 + 1)^4} &&\\
                &= \dfrac{-16x^4 + 32x^2 - 16 + 64x^4 - 64x^2}{(-x^2 + 1)^4} &&\\
                &= \dfrac{48x^4 - 32x^2 - 16}{(-x^2 + 1)^4} &&\\
                &= \dfrac{16(3x^4 - 2x^2 - 1)}{(-x^2 + 1)^4} &&\\
                &= \dfrac{16(3x^4 - 3x^2 + x^2 - 1)}{(-x^2 + 1)^4} &&\\
                &= \dfrac{16(3x^2(x^2 - 1) + (x^2 - 1))}{(-x^2 + 1)^4} &&\\
                &= \dfrac{16(x^2 - 1)(3x^2 + 1)}{(-x^2 + 1)^4} &&\\
                &= \dfrac{-16(3x^2 + 1)}{(-x^2 + 1)^3}
            \end{flalign*}
        \end{minipage}
        \item Zeichnung: \\
        \raisebox{-\height}{
            \begin{tikzpicture}
                % Set variables
                \pgfmathsetmacro{\axisPosWidth}{7};
                \pgfmathsetmacro{\axisNegWidth}{-7};
                \pgfmathsetmacro{\axisPosHeight}{7};
                \pgfmathsetmacro{\axisNegHeight}{-7};
                \pgfmathsetmacro{\xIncrement}{1.0};
                \pgfmathsetmacro{\yIncrement}{1.0};
                \pgfmathsetmacro{\xScale}{2};
                \pgfmathsetmacro{\yScale}{2};
                \pgfmathsetmacro{\gridxIncrement}{0.5};
                \pgfmathsetmacro{\gridyIncrement}{0.5};

                % Axis/grid limits
                \pgfmathsetmacro{\xUpperLimitAxis}{\xIncrement * floor(\axisPosWidth / \xIncrement)}
                \pgfmathsetmacro{\xLowerLimitAxis}{\xIncrement * ceil(\axisNegWidth / \xIncrement)}
                \pgfmathsetmacro{\yUpperLimitAxis}{\yIncrement * floor(\axisPosHeight /\yIncrement)}
                \pgfmathsetmacro{\yLowerLimitAxis}{\yIncrement * ceil(\axisNegHeight / \yIncrement)}

                \pgfmathsetmacro{\xUpperLimitGrid}{\gridxIncrement * floor(\axisPosWidth / \gridxIncrement)}
                \pgfmathsetmacro{\xLowerLimitGrid}{\gridxIncrement * ceil(\axisNegWidth / \gridxIncrement)}
                \pgfmathsetmacro{\yUpperLimitGrid}{\gridyIncrement * floor(\axisPosHeight / \gridyIncrement)}
                \pgfmathsetmacro{\yLowerLimitGrid}{\gridyIncrement * ceil(\axisNegHeight / \gridyIncrement)}

                % Axis environment (PGFPlots)
                \begin{axis}[
                    xmin=\xLowerLimitGrid, xmax=\xUpperLimitGrid,
                    ymin=\yLowerLimitGrid, ymax=\yUpperLimitGrid,
                    axis lines=middle,
                    axis line style={->, >=stealth},
                    grid=both,
                    xlabel={$x$},
                    ylabel={$f(x)$},
                    xlabel style={at={(axis description cs:1,0.55)}, anchor=north},
                    ylabel style={at={(axis description cs:0.55,1)}, anchor=south},
                    xtick={\xLowerLimitAxis,\xLowerLimitAxis+1,...,\xUpperLimitAxis},
                    ytick={\yLowerLimitAxis,\yLowerLimitAxis+1,...,\yUpperLimitAxis},
                    xticklabel={\pgfmathparse{\tick*\xScale}\pgfmathprintnumber{\pgfmathresult}},
                    yticklabel={\pgfmathparse{\tick*\yScale}\pgfmathprintnumber{\pgfmathresult}},
                    tick label style={font=\scriptsize},
                    x=1cm,
                    y=1cm,
                    minor tick num=1,
                    scale only axis,
                    grid=both,
                    restrict y to domain=\yLowerLimitGrid:\yUpperLimitGrid,
                ]

                % Example rational function
                \addplot[thick, myb, samples=1000, domain=\xLowerLimitGrid:0.96*\xUpperLimitGrid]
                    {(-(\x*\xScale)^2 + 9)/(-(\x*\xScale)^2 + 1)*(1/\yScale)};
                \node[anchor=north] at (axis cs:{2*(1/\xScale)},{-3*(1/\yScale)}) {\textcolor{myb}{$G_f$}};

                % v_1(x)
                \addplot[thick, myorange] coordinates {(-1*(1/\xScale), 1*\yLowerLimitGrid) (-1*(1/\xScale), 0.96*\yUpperLimitGrid)};
                \node[anchor=east] at (axis cs:{-1.25*(1/\xScale)},{12*(1/\yScale)}) {\textcolor{myorange}{$v_1(x)$}};

                % v_2(x)
                \addplot[thick, myorange] coordinates {(1*(1/\xScale), 1*\yLowerLimitGrid) (1*(1/\xScale),0.96*\yUpperLimitGrid)};
                \node[anchor=west] at (axis cs:{1.25*(1/\xScale)},{12*(1/\yScale)}) {\textcolor{myorange}{$v_2(x)$}};

                % h_1(x)
                \addplot[thick, myorange, domain=\xLowerLimitGrid:0.96*\xUpperLimitGrid]
                    {1*(1/\yScale)};
                \node[anchor=south] at (axis cs:{3*(1/\xScale)},{1.20*(1/\yScale)}) {\textcolor{myorange}{$h_1(x)$}};

                % Nullstellen
                \fill[black] (axis cs:{-3*(1/\xScale)},{0*(1/\yScale)}) circle[radius=2pt]; 
                \node[anchor=north] at (axis cs:{-3*(1/\xScale)},{-0.1*(1/\yScale)}) {\textcolor{black}{$N_1$}};

                \fill[black] (axis cs:{3*(1/\xScale)},{0*(1/\yScale)}) circle[radius=2pt]; 
                \node[anchor=north] at (axis cs:{3*(1/\xScale)},{-0.1*(1/\yScale)}) {\textcolor{black}{$N_2$}};

                \end{axis}
            \end{tikzpicture}
        }\\[3.0ex]
    \end{itemize}

\end{solution}
